% !TeX root = ../main.tex
\chapter{TRIỂN KHAI VÀ CẤU HÌNH HỆ THỐNG}
\label{ch:trienkhai}

\section{Môi trường và công cụ}
\label{sec:moitruong}

Hệ thống được triển khai thử nghiệm trên máy cá nhân chạy hệ điều hành Linux,
với các công cụ chính được liệt kê trong Bảng~\ref{tab:tools}.

\begin{table}[htbp]
\centering
\caption{Các công cụ và công nghệ sử dụng}
\label{tab:tools}
\small
\begin{tabular}{p{3.2cm} p{3.2cm} p{6.4cm}}
\toprule
\textbf{Thành phần} & \textbf{Công nghệ} & \textbf{Vai trò}\\
\midrule
Nhà cung cấp định danh & Keycloak 26 & Cung cấp OIDC/OAuth2, SSO, MFA\\
Cơ sở dữ liệu & PostgreSQL 16 & Lưu trữ dữ liệu người dùng, realm\\
Container & Docker & Đóng gói và chạy dịch vụ\\
Orchestration & k3s & Cụm Kubernetes nhẹ\\
Quản lý gói K8s & Helm & Triển khai Keycloak theo chart\\
CI/CD & GitHub Actions & Tự động hóa kiểm thử và build\\
GitOps & ArgoCD & Đồng bộ trạng thái từ Git\\
Giám sát & Prometheus/Grafana & Thu thập và trực quan hóa chỉ số\\
\bottomrule
\end{tabular}
\end{table}

\section{Triển khai Keycloak và PostgreSQL bằng Docker Compose}
\label{sec:compose}

Ở giai đoạn đầu, Keycloak và PostgreSQL được triển khai cục bộ bằng Docker
Compose nhằm làm quen và kiểm thử luồng xác thực trước khi đưa lên Kubernetes.
Mã nguồn cấu hình tham khảo được trình bày trong
Đoạn mã~\ref{lst:compose}.

\begin{lstlisting}[caption={Cấu hình Docker Compose cho Keycloak và PostgreSQL},label={lst:compose}]
services:
  postgres:
    image: postgres:16
    environment:
      POSTGRES_DB: keycloak
      POSTGRES_USER: keycloak
      POSTGRES_PASSWORD: ${DB_PASSWORD}
    volumes:
      - pgdata:/var/lib/postgresql/data
    healthcheck:
      test: ["CMD-SHELL", "pg_isready -U keycloak"]
      interval: 10s
      retries: 5

  keycloak:
    image: quay.io/keycloak/keycloak:26.0
    command: start-dev --import-realm
    environment:
      KC_DB: postgres
      KC_DB_URL: jdbc:postgresql://postgres:5432/keycloak
      KC_DB_USERNAME: keycloak
      KC_DB_PASSWORD: ${DB_PASSWORD}
      KC_HOSTNAME: localhost
      KC_METRICS_ENABLED: "true"
      KC_HEALTH_ENABLED: "true"
      KEYCLOAK_ADMIN: admin
      KEYCLOAK_ADMIN_PASSWORD: ${KC_ADMIN_PASSWORD}
    ports:
      - "8080:8080"
    depends_on:
      postgres:
        condition: service_healthy

volumes:
  pgdata:
\end{lstlisting}

Việc tách cấu hình nhạy cảm (mật khẩu) ra biến môi trường giúp tránh đưa thông
tin bí mật vào mã nguồn -- một yêu cầu quan trọng của thiết kế an toàn.

\subsection{Kiểm tra và xác thực dịch vụ}
\label{subsec:kiemtra}

Sau khi khởi chạy, cần kiểm tra dịch vụ hoạt động đúng trước khi cấu hình ứng
dụng. Đoạn mã~\ref{lst:check} minh họa việc đọc tài liệu khám phá OIDC và xin
token bằng luồng client credentials.

\begin{lstlisting}[caption={Kiểm tra điểm cuối OIDC của Keycloak},label={lst:check}]
# Doc tai lieu discovery
curl -s http://localhost:8080/realms/demo/.well-known/openid-configuration | jq .

# Xin token bang client credentials
curl -s -X POST \
  http://localhost:8080/realms/demo/protocol/openid-connect/token \
  -d grant_type=client_credentials \
  -d client_id=demo-app \
  -d client_secret=$CLIENT_SECRET | jq .
\end{lstlisting}

\section{Cấu hình định danh như mã nguồn}
\label{sec:realm}

Cấu hình realm, client và role được mô tả trong một tệp JSON và nạp tự động
khi Keycloak khởi động (tham số \texttt{--import-realm}). Một phần nội dung
được minh họa trong Đoạn mã~\ref{lst:realm}.

\begin{lstlisting}[caption={Trích tệp cấu hình realm dưới dạng mã nguồn},label={lst:realm}]
{
  "realm": "demo",
  "enabled": true,
  "sslRequired": "external",
  "accessTokenLifespan": 300,
  "otpPolicyType": "totp",
  "clients": [
    {
      "clientId": "demo-app",
      "protocol": "openid-connect",
      "publicClient": false,
      "standardFlowEnabled": true,
      "redirectUris": ["http://localhost:3000/callback"]
    }
  ],
  "roles": {
    "realm": [
      { "name": "user" },
      { "name": "admin" }
    ]
  }
}
\end{lstlisting}

Trong giai đoạn nâng cao, cấu hình này được quản lý bằng
\texttt{keycloak-config-cli}~\cite{kc-config-cli} hoặc Keycloak Terraform
Provider~\cite{tf-keycloak} để có thể áp dụng lặp lại (idempotent) và kiểm soát
phiên bản chặt chẽ hơn.

\subsection{Áp dụng cấu hình tự động với keycloak-config-cli}
\label{subsec:configcli}

Để cấu hình được áp dụng lặp lại và tự động (idempotent), công cụ
\texttt{keycloak-config-cli} được sử dụng. Công cụ này đọc các tệp JSON từ một
thư mục và đồng bộ chúng lên Keycloak, phù hợp để tích hợp vào pipeline CI.
Đoạn mã~\ref{lst:configcli} minh họa cách gọi công cụ.

\begin{lstlisting}[caption={Áp dụng cấu hình realm bằng keycloak-config-cli},label={lst:configcli}]
docker run --rm \
  -e KEYCLOAK_URL=http://keycloak:8080 \
  -e KEYCLOAK_USER=admin \
  -e KEYCLOAK_PASSWORD=$KC_ADMIN_PASSWORD \
  -v $(pwd)/realm:/config \
  quay.io/adorsys/keycloak-config-cli:latest \
  --import.files.locations=/config/demo-realm.json
\end{lstlisting}

Nhờ cách tiếp cận này, mọi thay đổi cấu hình đều đi qua Git, có thể xem lại
lịch sử và khôi phục khi cần, thay vì thao tác thủ công trên giao diện quản trị.

\section{Ứng dụng demo và đăng nhập SSO}
\label{sec:appdemo}

Ứng dụng demo đóng vai trò Relying Party, sử dụng thư viện OIDC để thực hiện
luồng Authorization Code kèm PKCE với Keycloak. Sau khi đăng nhập thành công,
ứng dụng nhận ID Token, xác thực chữ ký token và hiển thị thông tin người dùng
(họ tên, email, vai trò). Đây là kịch bản minh họa trực quan cho tính năng SSO
của hệ thống.

% TODO: Chèn ảnh chụp màn hình (1) trang đăng nhập Keycloak,
%       (2) màn hình ứng dụng demo sau khi đăng nhập.

\subsection{Mã nguồn ứng dụng demo}
\label{subsec:appcode}

Ứng dụng demo được viết bằng Node.js với thư viện \texttt{openid-client}, thực
hiện luồng Authorization Code kèm PKCE. Đoạn mã~\ref{lst:app} trình bày phần
xử lý đăng nhập và nhận token.

\begin{lstlisting}[caption={Xu ly dang nhap OIDC trong ung dung demo (Node.js)},label={lst:app}]
const express = require('express');
const session = require('express-session');
const { Issuer, generators } = require('openid-client');

const app = express();
app.use(session({ secret: process.env.SESSION_SECRET,
  resave: false, saveUninitialized: false }));

let client;
async function init() {
  const issuer = await Issuer.discover('http://localhost:8080/realms/demo');
  client = new issuer.Client({
    client_id: 'demo-app',
    client_secret: process.env.CLIENT_SECRET,
    redirect_uris: ['http://localhost:3000/callback'],
    response_types: ['code']
  });
}

app.get('/login', (req, res) => {
  const code_verifier = generators.codeVerifier();
  const code_challenge = generators.codeChallenge(code_verifier);
  req.session.code_verifier = code_verifier;
  const url = client.authorizationUrl({
    scope: 'openid profile email',
    code_challenge,
    code_challenge_method: 'S256'
  });
  res.redirect(url);
});

app.get('/callback', async (req, res) => {
  const params = client.callbackParams(req);
  const tokenSet = await client.callback(
    'http://localhost:3000/callback', params,
    { code_verifier: req.session.code_verifier });
  req.session.user = tokenSet.claims();
  res.redirect('/');
});

init().then(() => app.listen(3000));
\end{lstlisting}

\section{Đóng gói và triển khai trên Kubernetes}
\label{sec:k8s-deploy}

Sau khi kiểm thử cục bộ, hệ thống được thiết kế để triển khai trên cụm k3s.
Keycloak được triển khai thông qua Helm chart với các giá trị tùy biến
(Đoạn mã~\ref{lst:helm}), cho phép cấu hình số bản sao, tài nguyên và kết nối
cơ sở dữ liệu.

\begin{lstlisting}[caption={Trích tệp values.yaml cho Helm chart Keycloak},label={lst:helm}]
replicas: 2
image:
  repository: quay.io/keycloak/keycloak
  tag: "26.0"
command:
  - start
  - --optimized
extraEnv: |
  - name: KC_DB
    value: postgres
  - name: KC_HEALTH_ENABLED
    value: "true"
  - name: KC_METRICS_ENABLED
    value: "true"
resources:
  requests:
    cpu: "250m"
    memory: "512Mi"
  limits:
    cpu: "1"
    memory: "1Gi"
\end{lstlisting}

\subsection{Manifest triển khai Keycloak}
\label{subsec:manifest}

Đoạn mã~\ref{lst:k8sdeploy} và~\ref{lst:k8ssvc} trình bày manifest Deployment
và Service cho Keycloak, kèm các đầu dò sức khỏe để Kubernetes tự phát hiện Pod
lỗi.

\begin{lstlisting}[caption={Manifest Deployment cho Keycloak},label={lst:k8sdeploy}]
apiVersion: apps/v1
kind: Deployment
metadata:
  name: keycloak
  namespace: identity
spec:
  replicas: 2
  selector:
    matchLabels:
      app: keycloak
  template:
    metadata:
      labels:
        app: keycloak
    spec:
      containers:
        - name: keycloak
          image: quay.io/keycloak/keycloak:26.0
          args: ["start", "--optimized"]
          envFrom:
            - secretRef:
                name: keycloak-secret
          ports:
            - containerPort: 8080
          readinessProbe:
            httpGet: { path: /health/ready, port: 8080 }
          livenessProbe:
            httpGet: { path: /health/live, port: 8080 }
\end{lstlisting}

\begin{lstlisting}[caption={Service, Ingress và HorizontalPodAutoscaler},label={lst:k8ssvc}]
apiVersion: v1
kind: Service
metadata:
  name: keycloak
  namespace: identity
spec:
  selector:
    app: keycloak
  ports:
    - port: 8080
      targetPort: 8080
---
apiVersion: autoscaling/v2
kind: HorizontalPodAutoscaler
metadata:
  name: keycloak
  namespace: identity
spec:
  scaleTargetRef:
    apiVersion: apps/v1
    kind: Deployment
    name: keycloak
  minReplicas: 2
  maxReplicas: 5
  metrics:
    - type: Resource
      resource:
        name: cpu
        target:
          type: Utilization
          averageUtilization: 70
\end{lstlisting}

\section{Tự động hóa bằng CI/CD và GitOps}
\label{sec:cicd}

Quy trình CI/CD được xây dựng bằng GitHub Actions. Mỗi lần đẩy mã nguồn, hệ
thống tự động kiểm thử, xây dựng image ứng dụng và đẩy lên registry
(Đoạn mã~\ref{lst:ci}).

\begin{lstlisting}[caption={Trích quy trình CI bằng GitHub Actions},label={lst:ci}]
name: ci
on:
  push:
    branches: [ main ]
jobs:
  build:
    runs-on: ubuntu-latest
    steps:
      - uses: actions/checkout@v4
      - name: Build image
        run: docker build -t ghcr.io/${{ github.repository }}:${{ github.sha }} .
      - name: Scan image
        run: trivy image ghcr.io/${{ github.repository }}:${{ github.sha }}
      - name: Push image
        run: docker push ghcr.io/${{ github.repository }}:${{ github.sha }}
\end{lstlisting}

Về phía triển khai, ArgoCD~\cite{argocd} theo dõi kho Git chứa manifest và tự
động đồng bộ lên cụm k3s (Đoạn mã~\ref{lst:argo}). Nhờ vậy, trạng thái của hệ
thống luôn phản ánh đúng cấu hình trong Git.

\begin{lstlisting}[caption={Khai báo ứng dụng ArgoCD},label={lst:argo}]
apiVersion: argoproj.io/v1alpha1
kind: Application
metadata:
  name: keycloak
  namespace: argocd
spec:
  project: default
  source:
    repoURL: https://github.com/<user>/<repo>.git
    path: k8s/keycloak
    targetRevision: main
  destination:
    server: https://kubernetes.default.svc
    namespace: identity
  syncPolicy:
    automated:
      prune: true
      selfHeal: true
\end{lstlisting}

\section{Giám sát và vận hành}
\label{sec:giamsat}

Để vận hành dịch vụ định danh một cách tin cậy, cần thu thập chỉ số và log.
Keycloak hỗ trợ xuất chỉ số qua điểm cuối \texttt{/metrics} khi được bật. Cấu
hình Prometheus thu thập các chỉ số này được minh họa trong
Đoạn mã~\ref{lst:prom}.

\begin{lstlisting}[caption={Cấu hình Prometheus thu thập chỉ số Keycloak},label={lst:prom}]
scrape_configs:
  - job_name: keycloak
    metrics_path: /metrics
    static_configs:
      - targets: ["keycloak.identity.svc:8080"]
\end{lstlisting}

Trên Grafana, các bảng theo dõi chính gồm: số lần đăng nhập thành công/thất bại,
tỉ lệ cấp token, độ trễ xử lý yêu cầu, mức sử dụng CPU/bộ nhớ của các Pod và số
lượng phiên đang hoạt động. Việc giám sát giúp phát hiện sớm bất thường như tấn
công dò mật khẩu hoặc quá tải dịch vụ.

\section{Quản lý TLS với cert-manager}
\label{sec:tls}

Để đảm bảo mọi kết nối đều được mã hóa, chứng chỉ TLS được cấp và gia hạn tự
động thông qua cert-manager kết hợp với Let's Encrypt. Đoạn mã~\ref{lst:issuer}
trình bày khai báo một ClusterIssuer.

\begin{lstlisting}[caption={Khai bao ClusterIssuer cho Let's Encrypt},label={lst:issuer}]
apiVersion: cert-manager.io/v1
kind: ClusterIssuer
metadata:
  name: letsencrypt-prod
spec:
  acme:
    server: https://acme-v02.api.letsencrypt.org/directory
    email: admin@example.com
    privateKeySecretRef:
      name: letsencrypt-prod
    solvers:
      - http01:
          ingress:
            class: traefik
\end{lstlisting}

Ingress của Keycloak và ứng dụng demo chỉ cần khai báo annotation tương ứng, hệ
thống sẽ tự động lấy và gia hạn chứng chỉ mà không cần thao tác thủ công.

\section{Quản lý bí mật và cấu trúc kho GitOps}
\label{sec:gitops-secret}

Để có thể lưu bí mật an toàn trong Git, giải pháp Sealed Secrets được sử dụng:
bí mật được mã hóa bằng khóa công khai của cụm và chỉ có thể giải mã bởi
controller trong cụm. Đoạn mã~\ref{lst:sealed} minh họa một SealedSecret.

\begin{lstlisting}[caption={Vi du mot SealedSecret},label={lst:sealed}]
apiVersion: bitnami.com/v1alpha1
kind: SealedSecret
metadata:
  name: keycloak-secret
  namespace: identity
spec:
  encryptedData:
    KC_DB_PASSWORD: AgBy3i4OJSWK+PiTySY...
\end{lstlisting}

Kho GitOps được tổ chức theo mô hình base/overlay để tái sử dụng cấu hình giữa
các môi trường, như minh họa trong Đoạn mã~\ref{lst:gitrepo}.

\begin{lstlisting}[caption={Cau truc kho GitOps},label={lst:gitrepo}]
gitops/
|-- apps/
|   |-- keycloak/
|   `-- demo-app/
|-- base/
|   `-- keycloak/
`-- overlays/
    |-- dev/
    `-- prod/
\end{lstlisting}

\section{Quy trình triển khai từng bước}
\label{sec:quytrinh}

Quy trình triển khai từ mã nguồn tới môi trường chạy gồm các bước: (1) lập
trình viên đẩy thay đổi lên Git; (2) CI chạy kiểm thử và build image; (3) image
được đẩy lên registry; (4) manifest được cập nhật; (5) ArgoCD phát hiện và đồng
bộ lên cụm; (6) Kubernetes thực hiện cuốn chiếu; (7) hệ thống giám sát ghi nhận
trạng thái. Mỗi bước đều có thể kiểm tra và quay lui khi cần.

\section{Kiểm thử khả năng tự phục hồi}
\label{sec:recovery}

Để kiểm chứng khả năng tự phục hồi, có thể xóa một Pod của Keycloak và quan sát
Kubernetes tự tạo Pod mới. Bảng~\ref{tab:mttr} ghi nhận kết quả kiểm thử theo
chỉ số thời gian phục hồi trung bình (MTTR), đồng thời kiểm tra dịch vụ không
gián đoạn nhờ các bản sao còn lại.

\begin{table}[htbp]
\centering
\caption{Kết quả kiểm thử tự phục hồi (khung ghi nhận)}
\label{tab:mttr}
\small
\begin{tabular}{p{5.4cm} p{3.2cm} p{3.8cm}}
\toprule
\textbf{Kịch bản} & \textbf{MTTR} & \textbf{Gián đoạn}\\
\midrule
Xóa một Pod Keycloak & (đo) & Không\\
Xóa Pod cơ sở dữ liệu & (đo) & Ngắn\\
Cuốn chiếu phiên bản mới & (đo) & Không\\
\bottomrule
\end{tabular}
\end{table}

\section{Triển khai cơ sở dữ liệu và bảo mật mạng}
\label{sec:postgres-k8s}

Cơ sở dữ liệu PostgreSQL được triển khai bằng StatefulSet để đảm bảo dữ liệu
được lưu trên vùng nhớ ổn định (Đoạn mã~\ref{lst:pg}). Song song đó, chính sách
mạng (NetworkPolicy) được áp dụng để chỉ cho phép Keycloak kết nối tới cơ sở dữ
liệu, giảm bề mặt tấn công (Đoạn mã~\ref{lst:netpol}).

\begin{lstlisting}[caption={Manifest StatefulSet cho PostgreSQL},label={lst:pg}]
apiVersion: apps/v1
kind: StatefulSet
metadata:
  name: postgres
  namespace: identity
spec:
  serviceName: postgres
  replicas: 1
  selector:
    matchLabels: { app: postgres }
  template:
    metadata:
      labels: { app: postgres }
    spec:
      containers:
        - name: postgres
          image: postgres:16
          ports:
            - containerPort: 5432
          envFrom:
            - secretRef: { name: postgres-secret }
          volumeMounts:
            - name: data
              mountPath: /var/lib/postgresql/data
  volumeClaimTemplates:
    - metadata: { name: data }
      spec:
        accessModes: ["ReadWriteOnce"]
        resources:
          requests: { storage: 5Gi }
\end{lstlisting}

\begin{lstlisting}[caption={Chinh sach mang gioi han truy cap co so du lieu},label={lst:netpol}]
apiVersion: networking.k8s.io/v1
kind: NetworkPolicy
metadata:
  name: allow-keycloak-to-postgres
  namespace: identity
spec:
  podSelector:
    matchLabels: { app: postgres }
  ingress:
    - from:
        - podSelector:
            matchLabels: { app: keycloak }
      ports:
        - port: 5432
\end{lstlisting}

\section{Kết chương}

Chương này đã trình bày quá trình triển khai hệ thống từ môi trường cục bộ với
Docker Compose, cấu hình định danh như mã nguồn, ứng dụng demo SSO, cho tới
thiết kế triển khai trên Kubernetes cùng quy trình CI/CD và GitOps. Chương tiếp
theo đánh giá kết quả đạt được và phân tích các khía cạnh bảo mật.
